\subsection{Метод потенциалов}

Для полученного начального приближения проверяется, что число заполненных клеток
равно $m + n - 1$.

\begin{idea}
Это условие невырожденности опорного плана. Ограничений в закрытой задаче
$m + n$, но они линейно зависимы — сумма «строчных» равенств совпадает с суммой
«столбцовых» (и то, и другое равно $\sum a_i = \sum b_j$). Значит независимых
ограничений $m + n - 1$, и опорный вектор имеет ровно столько положительных
компонент (см. определение~\ref{def:support-vector}). Если заполненных клеток
меньше, план вырожден, и в одну из пустых клеток вписывают «нулевую перевозку»,
чтобы восстановить базис.

В примере~\ref{ex:transport} $m + n - 1 = 3 + 4 - 1 = 6$ — ровно столько клеток
и заполнено обоими методами.
\end{idea}

\subsubsection*{Алгоритм}

\paragraph{1) Проверка оптимальности решения.}
Сопоставим каждому столбцу переменную $v_j$, а каждой строке — переменную $u_i$;
эти переменные называются \textbf{потенциалами}. Для \emph{заполненных} клеток
выписываем равенства $v_j - u_i = c_{ij}$.

Для плана, построенного методом северо-западного угла в примере~\ref{ex:transport}:
\[
  \begin{cases}
    v_1 - u_1 = 14, \quad v_1 - u_2 = 10, \\
    v_2 - u_2 = 17, \quad v_3 - u_2 = 15, \\
    v_3 - u_3 = 25, \quad v_4 - u_3 = 21 .
  \end{cases}
\]
Уравнений $m + n - 1 = 6$, неизвестных $m + n = 7$, поэтому одну из них задают
произвольно. Положим $v_1 = 0$, тогда
\[
  u_1 = -14, \quad u_2 = -10, \quad u_3 = -20,
  \qquad v_2 = 7, \quad v_3 = 5, \quad v_4 = 1 .
\]

\paragraph{2) Проверка по пустым (зачёркнутым) клеткам.}
\[
  v_j - u_i \le c_{ij} \quad \forall i, j
  \qquad\Longrightarrow\qquad \text{план оптимальный} .
\]

\paragraph{3) Построение цикла пересчёта.}
Выбирая клетку, у которой $\abs{v_j - u_i - c_{ij}}$ максимальна, быстрее придём
к оптимальному решению, но в общем случае берём любую пустую клетку и строим
ломаную (\textbf{цикл пересчёта}) по принципу:
\begin{itemize}[nosep,left=0pt]
  \item направление только параллельно рёбрам клеток;
  \item изменение направления происходит только в заполненных клетках;
  \item ломаная всегда возвращается в клетку старта.
\end{itemize}
Такую ломаную можно построить всегда и только одним способом.

Вершины ломаной нумеруем, чередуя знаки $\oplus$ и $\ominus$, начиная с $\oplus$
в стартовой клетке. Выбираем все клетки со знаком $\ominus$ и объёмы перевозок,
им соответствующие:
\[
  \theta = \min_{\ominus} x_{ij} .
\]

\inputfig{figures/fig-cycle}

В примере стартовой выбрана клетка $(1,4)$. Цикл:
\[
  (1,4)^{\oplus} \to (1,1)^{\ominus} \to (2,1)^{\oplus} \to (2,3)^{\ominus}
  \to (3,3)^{\oplus} \to (3,4)^{\ominus} \to (1,4) ,
\]
\[
  \theta = \min\{27,\ 1,\ 17\} = 1 .
\]

\paragraph{4) Пересчёт.}
К клеткам со знаком $\oplus$ прибавляем $\theta$, из клеток со знаком $\ominus$
вычитаем:
\[
  \begin{array}{lll}
    (1,4)^{\oplus}: & 0 + \theta = 1, &\qquad
    (2,3)^{\ominus}: \ 1 - \theta = 0, \\
    (1,1)^{\ominus}: & 27 - \theta = 26, &\qquad
    (3,3)^{\oplus}: \ 26 + \theta = 27, \\
    (2,1)^{\oplus}: & 6 + \theta = 7, &\qquad
    (3,4)^{\ominus}: \ 17 - \theta = 16 .
  \end{array}
\]
Клетка $(2,3)$ обнулилась и становится пустой, вместо неё заполнена $(1,4)$ —
число заполненных клеток снова $6$. Так получаем новый опорный план.

\paragraph{5) Переход на пункт 1).}

\begin{idea}[title={Что здесь на самом деле произошло}]
Проверим пункт~2 для всех пустых клеток, $u = (-14, -10, -20)$,
$v = (0, 7, 5, 1)$. Величина $\Delta_{ij} = v_j - u_i - c_{ij}$:
\begin{center}
\begin{tabular}{lrrrrrr}
\toprule
клетка & $(1,2)$ & $(1,3)$ & $(1,4)$ & $(2,4)$ & $(3,1)$ & $(3,2)$ \\
\midrule
$\Delta_{ij}$ & $-7$ & $-2$ & $-13$ & $-13$ & $+6$ & $-3$ \\
\bottomrule
\end{tabular}
\end{center}
Условие $\Delta_{ij} \le 0$ нарушено только в клетке $(3,1)$, значит план не
оптимален, и улучшает его именно она. Выбранная на лекции клетка $(1,4)$ имеет
$\Delta_{14} = -13 < 0$, поэтому пересчёт по ней стоимость не уменьшает, а
\emph{увеличивает} на $\abs{\Delta_{14}}\cdot\theta = 13 \cdot 1 = 13$:
было $1681$, стало $1694$. Цикл в лекции построен, чтобы показать механику
пересчёта, а не чтобы улучшить план.

Правильный шаг — стартовая клетка $(3,1)$. Её цикл:
\[
  (3,1)^{\oplus} \to (3,3)^{\ominus} \to (2,3)^{\oplus} \to (2,1)^{\ominus}
  \to (3,1) , \qquad \theta = \min\{26,\ 6\} = 6 ,
\]
после пересчёта $x_{31} = 6$, $x_{33} = 20$, $x_{23} = 7$, $x_{21} = 0$, и
стоимость падает на $\Delta_{31}\cdot\theta = 6 \cdot 6 = 36$: с $1681$ до $1645$.
Общее правило: пересчёт по клетке с $\Delta_{ij}$ меняет стоимость на
$-\Delta_{ij}\,\theta$, поэтому выбирать надо клетки с $\Delta_{ij} > 0$, а среди
них — с наибольшим $\Delta_{ij}$.
\end{idea}
